\chapter{\textcolor{gray}{Potentiels membranaires}}
\section{Principes fondamentaux de l'électricité}
Potentiels de la membrane
\hspace*{1em}Les principaux ions du liquide extracellulaire sont le sodium ($Na^+$) et le chlorure ($Cl^-$), tandis que le liquide intracellulaire est riche en potassium ($K^+$) et en molécules non pénétrantes chargées négativement (phosphates, protéines).

\hspace*{1em}Les phénomènes électriques produits par la distribution de ces ions se manifestent au niveau de la membrane plasmique et sont essentiels à la communication entre cellules, notamment chez les neurones.

\begin{figure}[h]
	\centering
	\includegraphics[width=0.65\linewidth]{pictures/ions.png}
	\caption{La force d'attraction électrique}
\end{figure}

\hspace*{1em}Les charges opposées s'attirent, celles de même signe se repoussent. Lorsqu'elles sont séparées, elles créent une \textbf{différence de potentiel} (exprimée en millivolts), capable de générer un \textbf{courant électrique} si une voie est disponible. Le courant dépend de cette différence de potentiel et de la résistance du milieu, selon la loi d'Ohm :
\[
I = \frac{V}{R}
\]

	Les liquides biologiques, riches en ions, sont de bons conducteurs. En revanche, la membrane plasmique, composée de lipides, est un isolant qui sépare deux compartiments conducteurs : l'intérieur et l'extérieur de la cellule.
\section{Potentiels de la membrane}
\subsection{Le potentiel du repos}

\hspace*{1em}Toutes les cellules au repos ont une différence de potentiel à travers leur membrane plasmique, l'intérieur étant chargé négativement par rapport à l'extérieur. Ce potentiel est appelé le potentiel de membrane au repos. Par convention, le liquide extracellulaire sert de référence pour la tension, et la polarité du potentiel de membrane est exprimée en fonction de la charge excédentaire à l'intérieur de la cellule. Par exemple, si la différence de potentiel est de 70 mV et que l'intérieur est négatif, on dit que le potentiel est de -70 mV.

\begin{figure}[h]
	\centering
	\includegraphics[width=0.65\linewidth]{pictures/cell voltage.png}
	\caption{Mesure du potentiel de la membrane}
\end{figure}

\hspace*{1em}Le potentiel de membrane au repos varie de -25 à -100 mV selon le type de cellule, les neurones ayant généralement des valeurs entre -40 et -90 mV. Il existe en raison d'un léger excédent d'ions négatifs à l'intérieur de la cellule et d'ions positifs à l'extérieur. Ces charges excessives s'accumulent près de la membrane plasmique, tandis que le reste des fluides intra- et extracellulaires reste électriquement neutre.

\begin{figure}[h]
	\centering
	\includegraphics[width=0.65\linewidth]{pictures/ions and cell.png}
	\caption{Les charges opposées s'accumulant de part et d'autre de la membrane cellulaire.}
\end{figure}

\hspace*{1em}Les principales concentrations d'ions affectant le potentiel de membrane sont le sodium ($Na^+$), le potassium ($K^+$) et le chlorure ($Cl^-$). Le sodium et le potassium jouent un rôle essentiel dans la génération du potentiel de membrane au repos. La différence de concentration de ces ions est maintenue par la pompe sodium-potassium ($Na^+/K^+$-ATPase), qui expulse $Na^+$ et fait entrer $K^+$. Les concentrations de $Na^+$ et $Cl^-$ sont plus faibles à l'intérieur de la cellule, tandis que la concentration de $K^+$ est plus élevée à l'intérieur.

\begin{figure}[h]
	\centering
	\begin{minipage}{0.4\textwidth}
		\centering
		\includegraphics[width=\linewidth]{pictures/pomp 1.png}
	\end{minipage}
	\hspace{0.2cm}
	\begin{minipage}{0.4\textwidth}
		\centering
		\includegraphics[width=\linewidth]{pictures/pomp 2.png}
	\end{minipage}
	\caption{La pompe ATPase}
\end{figure}

\subsection{Potentiel d'action}

\hspace*{1em}Le potentiel d'action est un changement rapide et transitoire du potentiel membranaire dû à l'ouverture séquentielle de canaux voltage-dépendants pour $Na^+$ et $K^+$.

\vspace{1em}
\hspace*{1em}Les potentiels d'action diffèrent des autres types de potentiels (potentiels gradués) par leur amplitude et rapidité. Un potentiel d'action peut entraîner un changement de potentiel membranaire allant jusqu'à 100~mV, comme une dépolarisation de -70~mV à +30~mV, suivie d'une repolarisation. Ce processus est rapide (1--4 millisecondes) et peut se répéter plusieurs centaines de fois par seconde, permettant la communication du système nerveux sur de longues distances.
\begin{figure}[h]
	\centering
	\includegraphics[width=\linewidth]{pictures/2.png}
	\caption{Les variations du (a) potentiel de membrane (mV)}
\end{figure}

\textbf{Mécanisme du potentiel d'action :}\\
Au repos, la membrane est proche du \textbf{potentiel d'équilibre} du $K^+$. Une stimulation dépolarisante ouvre les canaux $Na^+$, entraînant une \textbf{dépolarisation} rapide via une boucle de rétroaction positive. La membrane devient temporairement positive.
\\
Ensuite, les canaux $Na^+$ se ferment via une \textit{inactivation}, et les canaux $K^+$ s'ouvrent, entraînant une \textbf{repolarisation}. Une \textbf{hyperpolarisation} temporaire suit, appelée \textbf{post-hyperpolarisation}, avant le \textbf{retour au repos}. Peu d'ions traversent la membrane à chaque potentiel d'action, mais les gradients sont maintenus par la pompe $Na^+/K^+$-ATPase.
\begin{figure}[h]
	\centering
	\includegraphics[scale=0.3]{pictures/1.png}
	\caption{Comportement des canaux $Na^+$ et $K^+$ voltage-dépendants. }
\end{figure}
\\
Des anesthésiques (ex. : lidocaïne) ou des toxines (ex. : tétrodotoxine du fugu) peuvent bloquer les canaux $Na^+$ et empêcher la conduction nerveuse.

Enfin, deux \textbf{périodes réfractaires} suivent chaque potentiel d'action : 

\begin{figure}[H]
	\centering
	\includegraphics[width=0.7\linewidth]{pictures/8.png}
	\caption{Périodes réfractaires absolue et relative du potentiel d'action}
\end{figure}
\begin{itemize}
	\item \textbf{Absolue} : aucun nouveau potentiel possible.
	\item \textbf{Relative} : un nouveau potentiel est possible si le stimulus est plus fort.\\
	Ces périodes assurent la séparation des signaux et leur propagation unidirectionnelle.
\end{itemize}          
\begin{figure}[H]
	\centering
	\includegraphics[width=0.65\linewidth]{pictures/3.png}
	\caption{Perméabilité relative de la membrane au sodium et au potassium durant un potentiel d'action}
\end{figure}

\begin{figure}[H]
	\centering
	\begin{minipage}{0.8\textwidth}
		\centering
		\includegraphics[scale=0.24]{pictures/6.png}
	\end{minipage}
	\begin{minipage}{0.8\textwidth}
		\centering
		\vspace{1cm}
		\includegraphics[scale=0.28]{pictures/7.png}
	\end{minipage}
	\caption{Relation entre le potentiel de membrane et l'intensité croissante d'un stimulus excitateur}
\end{figure}
Le potentiel d'action suit la règle du \textbf{tout ou rien} : s'il atteint le seuil (généralement $\sim$15~mV au-dessus du repos), il se déclenche entièrement, sinon, il échoue. L'intensité du stimulus ne change pas l'amplitude du potentiel d'action, mais affecte sa fréquence.
\begin{figure}[H]
	\centering
	\begin{minipage}{0.9\textwidth}
		\includegraphics[width=\linewidth]{pictures/4.png}
	\end{minipage}
	\begin{minipage}{0.9\textwidth}
		\centering
		\includegraphics[width=\linewidth]{pictures/5.png}
	\end{minipage}
	\caption{Mécanisme de rétroaction des canaux ioniques voltage-dépendants}
\end{figure}

\subsection{Transmission de l'information nerveuse : Synapses}

\begin{itemize}
	\item [$\diamond$] Les potentiels de membrane sont à la base de l'excitabilité des neurones, permettant la génération de signaux électriques.
	\item [$\to$] \textbf{La synapse} est le point de jonction où un neurone (présynaptique) transmet un signal à un autre neurone (postsynaptique) ou à une cellule cible, comme les muscles ou les glandes. Elle joue un rôle clé dans le traitement et la transmission de l'information dans le système nerveux.
	\item [$\diamond$] \textit{Structure d'une synapse chimique} :
	\begin{itemize}
		\item[-] Bouton présynaptique : Contient des vésicules remplies de neurotransmetteurs.
		\item[-] Fente synaptique : Espace étroit (20-40 nm) entre les deux neurones.
		\item[-] Membrane postsynaptique : Contient des récepteurs spécifiques aux neurotransmetteurs.
	\end{itemize}
\end{itemize}



\begin{figure}[h]
	\centering
	\begin{minipage}{0.9\textwidth}
		\includegraphics[width=0.9\linewidth]{pictures/synapse1}
	\end{minipage}
	\begin{minipage}{0.9\textwidth}
		\centering
		\includegraphics[width=\linewidth]{pictures/synapse2} 
	\end{minipage}
	\caption{Mécanismes d'une synapse chimique et détails de la libération des neurotransmetteurs}
\end{figure}

\textbf{Mécanisme de libération des neurotransmetteurs:}\\
\begin{itemize}
	\item [$\diamond$]\textit{Types de synapses} :
	\begin{itemize} 
		\item[-] Principalement \textbf{chimiques}, utilisant des neurotransmetteurs.
	\end{itemize}
	\item [$\diamond$] \textit{Mécanisme de transmission}:
	\begin{itemize}
		\item[-] Un potentiel d'action arrive au bouton présynaptique.
		\item[-] La dépolarisation ouvre des canaux calciques voltage-dépendants, permettant l'entrée de calcium.
		\item[-] Le calcium déclenche la libération des neurotransmetteurs dans la fente synaptique.
		\item[-] Les neurotransmetteurs se lient aux récepteurs postsynaptiques, modifiant le potentiel membranaire.
		\item[-] Les neurotransmetteurs sont ensuite dégradés (ex. : par des enzymes) ou recaptés pour arrêter le signal.
	\end{itemize}
\end{itemize}